\documentclass{beamer}

\usepackage{beamerthemeSingapore}
\usepackage[shorthands=off,brazil]{babel}
\usepackage[T1]{fontenc}
% sem guillemets: "> fora do código e << >> na fonte monoespaçada
\usepackage[protrusion=false,expansion=false]{microtype}
\DisableLigatures[<,>]{encoding = T1, family = tt*}
\usepackage[utf8]{inputenc}
\usepackage{graphicx}
\usepackage{url}
\usepackage{xcolor}
\usepackage{booktabs}
\usepackage{ragged2e}
\usepackage{etoolbox}
\usepackage{minted}
\usepackage{array}

\newcommand{\tcb}[1]{\textcolor{blue}{#1}}
\newcommand{\tcr}[1]{\textcolor{red}{#1}}
\newcommand{\tcg}[1]{\textcolor{green!50!black}{#1}}

\setminted{
  fontsize=\footnotesize,
  breaklines=true,
  breakanywhere=true,
  frame=lines,
  framesep=2mm,
  baselinestretch=1.05,
  tabsize=2
}

\apptocmd{\frame}{}{\justifying}{}

\setbeamertemplate{navigation symbols}{}

\title{12 - PHP: SuperGlobals, formulários, validação e includes}
\author{Prof. Alex Kutzke}
\date{Novembro de 2026}

\begin{document}

% ==========================================
\begin{frame}
	\begin{center}
		\large{Setor de Educação Profissional e Tecnológica}\\
		Tecnologia em Análise e Desenvolvimento de Sistemas\\
		\vspace{1cm}
		\small{DS122 - Desenvolvimento de Aplicações Web 1}
	\end{center}
	\titlepage
\end{frame}

% ==========================================
\begin{frame}{Objetivos da aula}
	Ao final desta aula você deve ser capaz de:
	\begin{enumerate}
		\item Ler dados da requisição com \texttt{\$\_GET}, \texttt{\$\_POST} e
		      \texttt{\$\_SERVER};
		\item Ligar um formulário HTML a um script PHP com \texttt{action},
		      \texttt{method} e \texttt{name};
		\item Escrever um formulário que envia os dados para a própria página;
		\item Validar dados no front-end e no back-end;
		\item Explicar XSS e SQL Injection, e como prevenir cada um;
		\item Dividir o código em arquivos com \texttt{include} e
		      \texttt{require}.
	\end{enumerate}
\end{frame}

% ==========================================
\begin{frame}{Roteiro}
	\tableofcontents
\end{frame}

% ==========================================
\begin{frame}{Código de exemplo}
	\begin{center}
		{\small\mbox{\texttt{https://gitlab.com/ds122-alexkutzke/ds122-php-forms-example}}}
	\end{center}

	\vspace{0.3cm}
	Bibliografia recomendada:
	\begin{itemize}
		\item Documentação oficial do PHP, Variáveis Pré-definidas:\\
		      {\small\mbox{\texttt{https://www.php.net/manual/pt\_BR/reserved.variables.php}}}
		\item w3schools, PHP SuperGlobals e PHP Forms.
	\end{itemize}
\end{frame}

% ==========================================
\section{SuperGlobals}
% ==========================================

\begin{frame}{Arrays SuperGlobals}
	Ao iniciar um script, o interpretador do PHP cria uma série de arrays
	associativos globais, as \textbf{SuperGlobals}.

	\vspace{0.2cm}
	\begin{itemize}
		\item Estão disponíveis em qualquer escopo do código, dentro ou fora de
		      funções;
		\item Trazem informações sobre o ambiente do servidor, os cabeçalhos HTTP
		      e, principalmente, os dados enviados pelo usuário na requisição.
	\end{itemize}
\end{frame}

\begin{frame}{As SuperGlobals mais utilizadas}
	\small
	\begin{tabular}{>{\ttfamily}l p{0.68\textwidth}}
		\toprule
		\$\_SERVER  & ambiente do servidor, cabeçalhos HTTP, caminhos e
		localização de scripts \\
		\$\_GET     & dados enviados pela URL (\emph{query string}) ou por
		formulário com método GET \\
		\$\_POST    & dados de formulário enviados com método POST, no corpo da
		requisição, e não na URL \\
		\$\_REQUEST & \texttt{\$\_GET} e \texttt{\$\_POST} combinados (e
		\texttt{\$\_COOKIE}, conforme a diretiva \texttt{request\_order}) \\
		\$\_FILES   & arquivos enviados por upload \\
		\$\_SESSION & dados da sessão do usuário \\
		\$\_COOKIE  & cookies enviados pelo navegador \\
		\$GLOBALS   & todas as variáveis globais do código \\
		\$\_ENV     & variáveis de ambiente \\
		\bottomrule
	\end{tabular}

	\vspace{0.2cm}
	Prefira \texttt{\$\_GET} ou \texttt{\$\_POST} a \texttt{\$\_REQUEST}, por
	segurança e clareza.
\end{frame}

\begin{frame}{\texttt{\$GLOBALS}}
	Dá acesso a todas as variáveis globais definidas no código.

	\vspace{0.3cm}
	Usá-lo para alterar, dentro de uma função, uma variável de fora dela é
	considerado má prática. Prefira passar parâmetros e usar o retorno da
	função.
\end{frame}

\begin{frame}[fragile]{Consulta ao \texttt{\$\_SERVER}: \texttt{server.php}}
	\begin{minted}[fontsize=\scriptsize]{php}
<?php
// Exibindo informações sobre o servidor e a requisição atual
echo "Caminho do script: " . htmlspecialchars($_SERVER['PHP_SELF']) . "<br>";
echo "Nome do servidor: " . $_SERVER['SERVER_NAME'] . "<br>";
echo "Host: " . $_SERVER['HTTP_HOST'] . "<br>";
echo "Navegador/User-Agent: " . $_SERVER['HTTP_USER_AGENT'] . "<br>";
?>
	\end{minted}

	Texto exibido, com \texttt{php -S localhost:8000} e a requisição feita pelo \texttt{curl}:

	\begin{minted}[fontsize=\scriptsize]{text}
Caminho do script: /server.php
Nome do servidor: localhost
Host: localhost:8000
Navegador/User-Agent: curl/8.18.0
	\end{minted}

	\texttt{htmlspecialchars()} no \texttt{PHP\_SELF} evita injeção de código,
	como veremos na seção de XSS.
\end{frame}

% ==========================================
\section{Formulários}
% ==========================================

\begin{frame}{Formulário: \texttt{action}, \texttt{method} e \texttt{name}}
	Para o formulário HTML enviar dados ao back-end, ele define dois atributos:

	\begin{enumerate}
		\item \textbf{\texttt{action}}: a URL (recurso ou arquivo PHP) para onde
		      os dados serão enviados;
		\item \textbf{\texttt{method}}: o verbo HTTP do envio, \texttt{GET} ou
		      \texttt{POST}.
	\end{enumerate}

	\vspace{0.3cm}
	Cada campo (\texttt{<input>}, \texttt{<select>}, \texttt{<textarea>})
	precisa do atributo \textbf{\texttt{name}}. O PHP usa o valor desse
	atributo como \textbf{chave} na SuperGlobal:

	\vspace{0.2cm}
	\begin{center}
		\texttt{<input name="termo\_busca">} \quad corresponde a \quad
		\texttt{\$\_GET["termo\_busca"]}
	\end{center}
\end{frame}

\begin{frame}[fragile]{Formulário de busca}
	\begin{minted}{html}
<!DOCTYPE html>
<html lang="pt-BR">
  <head>
    <meta charset="utf-8">
    <title>Teste de formulário</title>
  </head>
  <body>
    <form action="form_test.php" method="get">
      <label for="busca">Pesquisar:</label>
      <input type="text" id="busca" name="termo_busca">
      <button type="submit">Buscar</button>
    </form>
  </body>
</html>
	\end{minted}

	Ao clicar em ``Buscar'', o navegador faz a requisição para\\
	\texttt{form\_test.php?termo\_busca=valor\_digitado}.
\end{frame}

\begin{frame}[fragile]{\texttt{form\_test.php}}
	\begin{minted}[fontsize=\scriptsize]{html+php}
<?php
  // O operador ?? (Null Coalescing) verifica se o campo foi enviado.
  // Se não foi, atribui uma string vazia, evitando erros de "Undefined array key".
  $busca = $_GET["termo_busca"] ?? '';
?>
<!DOCTYPE html>
<html lang="pt-BR">
  <head>
    <meta charset="utf-8">
    <title>Resultado da Busca</title>
  </head>
  <body>
    <h2>Você buscou por:</h2>
    <!-- O htmlspecialchars protege contra a execução acidental de tags HTML digitadas pelo usuário -->
    <p><?= htmlspecialchars($busca) ?></p>
  </body>
</html>
	\end{minted}
\end{frame}

\begin{frame}{Formulário na mesma página (\emph{self-submitting})}
	O formulário e o código PHP que o processa ficam no mesmo arquivo.

	\vspace{0.2cm}
	\begin{itemize}
		\item O atributo \texttt{action} é omitido, ou aponta para o próprio
		      script;
		\item O script verifica o método da requisição para saber se deve só
		      exibir o formulário ou também processar os dados:
	\end{itemize}

	\vspace{0.2cm}
	\begin{center}
		\texttt{if (\$\_SERVER["REQUEST\_METHOD"] === "POST") \{ ... \}}
	\end{center}

	\vspace{0.2cm}
	O primeiro acesso à página é um GET; o envio do formulário, um POST.
\end{frame}

\begin{frame}[fragile]{Mesma página: o processamento}
	\begin{minted}{php}
<?php
$nome = '';
$mensagem = '';

// Verifica se o formulário foi submetido
if ($_SERVER["REQUEST_METHOD"] === "POST") {
    // Resgata o valor. Se não existir, assume string vazia.
    $nome = $_POST['fname'] ?? '';

    if (empty(trim($nome))) {
        $mensagem = "O nome é obrigatório!";
    } else {
        $mensagem = "Olá, " . htmlspecialchars($nome) . "!";
    }
}
?>
	\end{minted}
\end{frame}

\begin{frame}[fragile]{Mesma página: o formulário}
	Continuação do mesmo arquivo:

	\begin{minted}[fontsize=\scriptsize]{html+php}
<!DOCTYPE html>
<html>
<body>

<!-- htmlspecialchars impede que o PHP_SELF seja vetor de ataque -->
<form method="post" action="<?= htmlspecialchars($_SERVER['PHP_SELF']) ?>">
  <label for="fname">Nome:</label>
  <input type="text" id="fname" name="fname" value="<?= htmlspecialchars($nome) ?>">
  <button type="submit">Enviar</button>
</form>

<p><strong><?= $mensagem ?></strong></p>

</body>
</html>
	\end{minted}

	O \texttt{value} devolve ao campo o que foi digitado, e o usuário não
	precisa preencher de novo.
\end{frame}

% ==========================================
\section{Validação}
% ==========================================

\begin{frame}{Validação}
	Verificar se os dados recebidos seguem as regras de negócio esperadas e se
	são seguros. Nunca confie nos dados enviados pelo usuário.

	\vspace{0.3cm}
	A validação ocorre em duas camadas:
	\begin{enumerate}
		\item \textbf{Front-end} (\emph{client-side}): avisa o erro antes do
		      envio, e poupa processamento do servidor. Pode ser burlada
		      facilmente;
		\item \textbf{Back-end} (\emph{server-side}): a validação definitiva e
		      obrigatória. Garante a integridade e a segurança do sistema, porque
		      não pode ser manipulada pelo usuário.
	\end{enumerate}
\end{frame}

\begin{frame}[fragile]{Validação no front-end: atributos do HTML5}
	Os próprios atributos do HTML5 validam o formulário, sem JavaScript:

	\begin{itemize}
		\item \texttt{required}, \texttt{minlength}, \texttt{maxlength},
		      \texttt{min}, \texttt{max};
		\item \texttt{pattern}, com uma expressão regular;
		\item tipos semânticos, como \texttt{type="email"} e
		      \texttt{type="number"}.
	\end{itemize}

	\begin{minted}{html}
<input type="email" name="email" required>
<input type="text" name="nome" required minlength="3" maxlength="60">
<input type="number" name="quantidade" min="1" max="10">
	\end{minted}
\end{frame}

\begin{frame}[fragile]{Validação no front-end: JavaScript}
	Para lógicas mais complexas, o JavaScript intercepta o evento
	\texttt{submit} do formulário:

	\begin{minted}{html}
<form id="form-test" action="processa.php" method="post">
    <input type="text" name="nome" id="nome_input">
    <span id="erro-nome" style="color: red;"></span>
    <button type="submit">Enviar</button>
</form>
	\end{minted}
\end{frame}

\begin{frame}[fragile]{Validação no front-end: o \emph{script}}
	Continuação do mesmo arquivo:

	\begin{minted}[fontsize=\scriptsize]{html}
<script>
    document.getElementById("form-test").addEventListener("submit", function(event) {
        const nomeInput = document.getElementById("nome_input").value;
        const erroSpan = document.getElementById("erro-nome");

        // Se estiver vazio, exibe o erro e cancela o envio
        if (nomeInput.trim() === "") {
            erroSpan.textContent = "O nome é obrigatório.";
            event.preventDefault(); // Impede o formulário de ser enviado ao servidor
        } else {
            erroSpan.textContent = "";
        }
    });
</script>
	\end{minted}
\end{frame}

\begin{frame}[fragile]{Validação no back-end}
	O PHP verifica se os dados existem, se estão no formato correto, e os
	higieniza (sanitização). \texttt{trim()} remove os espaços nas extremidades:

	\begin{minted}[fontsize=\scriptsize]{php}
<?php
if ($_SERVER["REQUEST_METHOD"] === "POST") {
    $nome = trim($_POST["nome"] ?? ''); // Remove espaços no início e no final

    if (empty($nome)) {
        echo "O campo nome é inválido ou está vazio.";
    } else {
        // Converte caracteres como < e > para entidades HTML, na hora de exibir
        echo "Olá, " . htmlspecialchars($nome) . "!";
    }
}
?>
	\end{minted}

	\texttt{htmlspecialchars()} fica para a hora de exibir. Feita na entrada, a
	conversão gravaria no banco de dados \texttt{\&lt;} no lugar de
	\texttt{<}.
\end{frame}

\begin{frame}{\texttt{filter\_var}}
	Para validações avançadas, o PHP oferece a família de funções
	\textbf{\texttt{filter\_var}}, com regras prontas para e-mails, URLs e
	números.

	\vspace{0.3cm}
	Documentação:\\
	{\small\mbox{\texttt{https://www.php.net/manual/pt\_BR/function.filter-var.php}}}
\end{frame}

% ==========================================
\section{Segurança}
% ==========================================

\begin{frame}{Segurança}
	A internet é um ambiente hostil. Formulários são a principal porta de
	entrada para ataques contra aplicações Web.

	\vspace{0.3cm}
	Dois ataques:
	\begin{itemize}
		\item \textbf{Cross-site Scripting} (XSS);
		\item \textbf{SQL Injection} (injeção de SQL).
	\end{itemize}
\end{frame}

\begin{frame}[fragile]{Cross-site Scripting (XSS)}
	A aplicação recebe dados com scripts (como JavaScript) e os exibe para
	outros usuários sem tratamento.

	\vspace{0.2cm}
	Campo ``Nome'' preenchido com:

	\begin{minted}{html}
José <script>alert('Roubei seus cookies!');</script>
	\end{minted}

	Se a aplicação fizer apenas \texttt{echo \$\_POST['nome'];}, o navegador
	interpreta a tag \texttt{<script>} e executa o código.

	\vspace{0.2cm}
	\textbf{Prevenção}: use \textbf{sempre} \texttt{htmlspecialchars()} antes de
	exibir no HTML qualquer dado vindo do usuário. Ela transforma \texttt{<} em
	\texttt{\&lt;} e \texttt{>} em \texttt{\&gt;}, e o navegador trata o
	conteúdo apenas como texto.
\end{frame}

\begin{frame}[fragile]{SQL Injection}
	Trechos de SQL enviados pelo formulário alteram a consulta feita ao banco
	de dados.

	\begin{minted}{php}
$sql = "SELECT * FROM Usuarios WHERE id = " . $_POST['id'];
	\end{minted}

	O usuário digita no campo ID: \texttt{105 OR 1=1}. A consulta final fica:

	\begin{minted}{sql}
SELECT * FROM Usuarios WHERE id = 105 OR 1=1;
	\end{minted}

	Como \texttt{1=1} é sempre verdadeiro, o banco devolve os dados de
	\textbf{todos} os usuários.
\end{frame}

\begin{frame}{Prevenção de SQL Injection}
	\textbf{Prepared Statements} (consultas preparadas): o dado enviado pelo
	usuário é tratado separadamente da instrução SQL.

	\vspace{0.3cm}
	Existem tanto na extensão \textbf{mysqli}, que usaremos na aula de PHP e
	MySQL, quanto na \textbf{PDO}.
\end{frame}

% ==========================================
\section{Includes}
% ==========================================

\begin{frame}{Includes e requires}
	Sistemas grandes ficam divididos em arquivos menores: um para o cabeçalho,
	um para o rodapé, outro para a conexão com o banco.

	\vspace{0.2cm}
	\small
	\begin{tabular}{>{\ttfamily\footnotesize}l p{0.56\textwidth}}
		\toprule
		include 'arquivo.php'; & arquivo inexistente: \textbf{aviso}
		(\emph{Warning}), e a execução \textbf{continua} \\ \midrule
		require 'arquivo.php'; & arquivo inexistente: \textbf{erro fatal}
		(\emph{Fatal Error}), e a execução \textbf{para}. Mais seguro para
		arquivos essenciais \\ \midrule
		\begin{tabular}[t]{@{}l@{}}include\_once\\require\_once\end{tabular} & iguais aos pares, mas não incluem
		de novo um arquivo já incluído na mesma execução. Evitam redefinir
		funções e variáveis \\
		\bottomrule
	\end{tabular}
\end{frame}

\begin{frame}[fragile]{Exemplo de \texttt{require}}
	\begin{columns}[T]
		\begin{column}{0.27\textwidth}
			\texttt{variaveis.php}:
			\begin{minted}[fontsize=\scriptsize]{php}
<?php
$cor = 'verde';
$fruta = 'maçã';
?>
			\end{minted}
		\end{column}
		\begin{column}{0.71\textwidth}
			\texttt{teste.php}:
			\begin{minted}[fontsize=\scriptsize]{php}
<?php
// Tenta imprimir variáveis que ainda não existem
echo "Uma $fruta $cor. <br>";

// Inclui o arquivo que contém as variáveis
require 'variaveis.php';

// Agora as variáveis estão disponíveis no escopo
echo "Uma $fruta $cor."; // Saída: Uma maçã verde.
?>
			\end{minted}
		\end{column}
	\end{columns}

	\vspace{0.2cm}
	No primeiro \texttt{echo}, o PHP emite \texttt{Warning: Undefined variable
	\$fruta} (e o mesmo para \texttt{\$cor}) e imprime \verb|Uma  .|
\end{frame}

% ==========================================
\begin{frame}{Exercícios e próximos passos}
	Os exercícios desta aula não têm entrega e não valem nota. O enunciado
	está no fim do material desta aula, na seção Exercícios, e o conteúdo cai
	na Prova 3.

	\vspace{0.3cm}
	Próxima aula, em 13/11: PHP e MySQL. Conexão com \texttt{mysqli},
	\emph{prepared statements} e operações CRUD.

	\vspace{0.3cm}
	A Entrega 3 do trabalho (PHP e MySQL) vence em 27/11, às 23h59. Não há
	entrega em atraso: vale o último \emph{commit} enviado até o prazo.
\end{frame}

\end{document}
